\documentclass[a4paper]{article}

\usepackage{graphicx}
\usepackage{amsmath,amssymb}
\usepackage{minted}
\usepackage{multirow}
\usepackage{float}
\usepackage{algorithm}
\usepackage{algpseudocode}
\usepackage{todonotes}
\usepackage{forest}
\usepackage{xstring}
\usepackage{xcolor}
\usepackage[most]{tcolorbox}
\usepackage{xparse}
\usepackage{natbib}
\usepackage{adjustbox}

\usetikzlibrary{graphs,graphdrawing}
\usegdlibrary{layered}

% for external graphviz
\input{/home/donkarlo/Dropbox/repo/nd_language_ai_project/src/nd_language_ai/formal/computer/programming/tex/latex/code_clips/external_graphviz.tex}

% for tags
\input{/home/donkarlo/Dropbox/repo/nd_language_ai_project/src/nd_language_ai/formal/computer/programming/tex/latex/part/control_sequence/kind/semtag/semtag.tex}

% for links
\input{/home/donkarlo/Dropbox/repo/nd_language_ai_project/src/nd_language_ai/formal/computer/programming/tex/latex/code_clips/seehere.tex}

% for nested lists and sections
\input{/home/donkarlo/Dropbox/repo/nd_language_ai_project/src/nd_language_ai/formal/computer/programming/tex/latex/code_clips/deep_structure.tex}

\title{}
\date{}
\author{Mohammad Rahmani}

\begin{document}
   \maketitle
   \begin{table}[H]
      \centering
      \small
      \setlength{\tabcolsep}{4pt}
      \begin{tabular}{p{0.25\textwidth} p{0.18\textwidth} p{0.47\textwidth}}
         \hline
         \textbf{Concept} & \textbf{Variable name} & \textbf{Description} \\
         \hline

         Depolarization period
         & $T_{\mathrm{dep}}$
         & Duration of the depolarization period. \\

         Maximum learned certainty
         & $C_{\max}$
         & Maximum confidence that the state neuron can reach. \\

         Repolarization period
         & $T_{\mathrm{rep}}$
         & Duration of the repolarization period. \\

         Firing threshold
         & $\theta_{\mathrm{fire}}$
         & Evidence threshold above which the state neuron generates a spike. \\

         Refractory period
         & $T_{\mathrm{ref}}$
         & Duration after a spike during which another spike is not allowed. \\

         Refractory depth
         & $D_{\mathrm{ref}}$
         & Strength of the post-spike inhibition or reset. \\

         Resting membrane potential
         & $V_{\mathrm{rest}}$
         & Baseline evidence potential when the represented state is not active. \\

         \hline
      \end{tabular}
      \caption{Action-potential-related parameters of a state neuron.}
      \label{tab:state-neuron-action-potential-parameters}
   \end{table}

   \bibliographystyle{plainnat}
   \bibliography{/home/donkarlo/Dropbox/repo/nd_shared_research/bibliography/bibliography.bib}
\end{document}
